
\appendix
%==========================================================================
\section{Code source complet}\label{app:code}
%==========================================================================

Tout est également disponible sur \texttt{https://github.com/tomasnp/eals-spark}. Les
listings ci-dessous sont les fichiers tels quels (verbatim, en anglais dans le code
source).

\subsection{\texttt{src/spark\_utils.py} - session Spark}
\lstinputlisting{../src/spark_utils.py}

\subsection{\texttt{src/data.py} - chargement, filtrage 10-core, découpage}
\lstinputlisting{../src/data.py}

\subsection{\texttt{src/eals\_local.py} - référence NumPy et noyaux partagés}
\lstinputlisting{../src/eals_local.py}

\subsection{\texttt{src/eals\_rdd.py} - implémentation Spark RDD}
\lstinputlisting{../src/eals_rdd.py}

\subsection{\texttt{src/evaluate.py} - HR@$k$, NDCG@$k$, classement sur catalogue entier}
\lstinputlisting{../src/evaluate.py}

\subsection{\texttt{src/baselines.py} - ALS de MLlib et MostPopular}
\lstinputlisting{../src/baselines.py}

\subsection{\texttt{src/check\_correctness.py} - le contrôle d'exactitude}
\lstinputlisting{../src/check_correctness.py}

\subsection{\texttt{src/experiments.py} - la campagne de mesures}
\lstinputlisting{../src/experiments.py}

\subsection{\texttt{src/plots.py} - figures}
\lstinputlisting{../src/plots.py}

\subsection{\texttt{report/tables.py} - chaque chiffre de ce rapport}
\lstinputlisting{../report/tables.py}

\subsection{\texttt{scripts/env.sh}, \texttt{run\_all.sh}, \texttt{build\_report.sh}}
\lstinputlisting[language=bash]{../scripts/env.sh}
\lstinputlisting[language=bash]{../scripts/run_all.sh}
\lstinputlisting[language=bash]{../scripts/build_report.sh}

%==========================================================================
\section{Le notebook de démonstration}\label{app:nb}
%==========================================================================

\texttt{notebooks/demo\_small\_data.ipynb} exécute tout le pipeline de bout en bout
sur \textbf{ml-100k} (943 utilisateurs, 1\,152 items, 97\,953 interactions) en
quelques minutes sur un portable. Son but est la vérification, pas la mesure : un
lecteur devrait pouvoir exécuter « Run All » et voir, sans attendre, que chaque
affirmation de ce rapport est étayée par du code qui s'exécute. Chaque section
indique son entrée et sa sortie.

\begin{table}[h]
\centering\small
\caption{Cellules du notebook, avec entrée et sortie.}\label{tab:nb}
\begin{tabular}{@{}p{0.5cm}p{4.6cm}p{4.6cm}p{5.2cm}@{}}
\toprule
\S & Entrée & Sortie & Ce que cela démontre \\
\midrule
1 & le JDK installé & versions Java / Spark / NumPy, master Spark et parallélisme
  & l'environnement est bien celui décrit en Section~\ref{sec:setup} ; en particulier
    \texttt{JAVA\_HOME} pointe vers un JDK 17. \\
2 & \texttt{data/raw/ml-100k/u.data} (brut \texttt{user\,item\,rating\,ts}) & le
    dictionnaire de statistiques et les premières lignes du découpage d'entraînement
    & le filtrage 10-core, la ré-indexation dense et le découpage leave-one-out
    reproduisent le Tableau~\ref{tab:datasets}, ligne \texttt{ml-100k}. \\
3 & fréquences des items d'entraînement & $\min/\text{médiane}/\max$ de $c_i$ et
    $\sum_i c_i$ & l'Éq.~\eqref{eq:ci} est correctement implémentée :
    $\sum_i c_i=c_0$ pour tout $\alpha$, et $\alpha=0,4$ étale $c_i$ sur un rapport
    6:1 sur ml-100k alors que $\alpha=0$ le réduit à $c_0/N$. Cela montre aussi que
    l'item le plus populaire atteint $c_i>1$ - la réserve S2 de la
    Section~\ref{sec:discussion}. \\
4 & les paires d'entraînement sous forme de deux tableaux d'entiers & un tableau de
    $J$ par itération et sa décroissance & l'objectif~\eqref{eq:obj7} décroît de
    façon monotone : les règles de mise à jour~\eqref{eq:pu}--\eqref{eq:qi} sont
    correctes. \\
5 & le DataFrame d'entraînement, 8 partitions & temps réel par itération décomposé en
    étape utilisateur / étape item / broadcast / driver & l'anatomie d'une itération
    distribuée, et la part prise par la partie non parallèle. \\
6 & les modèles des \S4 et \S5 & $\max|\Delta P|$, $\max|\Delta Q|$ et une assertion
    (\texttt{assert}) & le contrôle d'exactitude de la
    Section~\ref{sec:correctness}, en miniature. La cellule \emph{échoue
    bruyamment} si les implémentations divergent. \\
7 & matrices de facteurs + l'ensemble de test leave-one-out & HR@10/@100 et
    NDCG@10/@100 pour MostPopular, eALS, eALS avec $\alpha=0$ et MLlib ALS & eALS bat
    le plancher de popularité avec une large marge, et $\alpha>0$ bat $\alpha=0$. \\
8 & les mêmes données réajustées avec 1, 2, 4 cœurs & secondes par itération,
    accélération, efficacité & une miniature de la Section~\ref{sec:strong}, et une
    honnête : sur ml-100k l'accélération ressort \emph{en dessous de un} - plus de
    cœurs rendent eALS plus lent, car les données sont bien trop petites pour Spark. \\
\bottomrule
\end{tabular}
\end{table}

\paragraph{Temps d'exécution.} Environ deux minutes de bout en bout sur la machine
de la Section~\ref{sec:setup}, dominées par les démarrages de session Spark
(\S8 redémarre la session une fois par nombre de cœurs, ce qui est inévitable :
\texttt{local[$p$]} est fixé à la création de la session).
